% v2-figures/delegation.tex: the delegation chain from a human to an agent to a service,
% as built at v9.354. Layout note: the chain is a single left-to-right band about 13cm
% wide and 3cm tall, with the two prose blocks BELOW it rather than beside it, so the
% figure's aspect ratio stays close to the text block and \figfile's adjustbox barely has
% to scale it. An earlier version ran the prose alongside the chain, which made the box
% twice as wide as tall and shrank every label past reading size.
%
% Colour carries meaning here, as everywhere in this paper: navy for what the HUMAN holds,
% accent for the signed artifacts that run, steel for the outside party, dashed gold for
% the one thing that is NOT closed.
\begin{tikzpicture}[node distance=0pt]
  % --- the chain -----------------------------------------------------------------
  \node[humanbox, minimum width=3.2cm, text width=2.9cm, minimum height=1.55cm] (h) at (0,0)
    {Человек\\[-1pt]{\mdseries\scriptsize удостоверение + ключ, привязанный эмитентом}};
  \node[corebox, minimum width=3.2cm, text width=2.9cm, minimum height=1.55cm] (g) at (5.0,0)
    {Полномочие с областью\\[-1pt]{\mdseries\scriptsize подписано ключом держателя}};
  \node[corebox, minimum width=3.2cm, text width=2.9cm, minimum height=1.55cm] (a) at (10.0,0)
    {Программный агент\\[-1pt]{\mdseries\scriptsize подписывает каждое действие}};
  \node[extbox, minimum width=3.2cm, text width=2.9cm, minimum height=1.55cm] (s) at (15.0,0)
    {Служба\\[-1pt]{\scriptsize проверяет автономно}};

  \draw[flow] (h) -- node[faint, above=1pt] {выдаёт} (g);
  \draw[flow] (g) -- node[faint, above=1pt] {предъявляет} (a);
  \draw[flow] (a) -- node[faint, above=1pt] {действует} (s);

  % what rides inside the grant, hung under it so the chain stays one clean band
  \node[lbl, below=5pt of g, text width=3.3cm] (gd)
    {действия · пределы · срок\\ дескриптор отзыва};
  \node[lbl, below=5pt of a, text width=2.9cm] (ad)
    {называет действие и собственное одноразовое число службы};

  % --- what the service is told --------------------------------------------------
  \node[anchor=north west, font=\scriptsize, text=navydark, text width=8.2cm, align=left]
    (checks) at (-1.6,-1.95) {
    \textbf{Что проверяет служба:}\\
    1. полномочие подлинно и не истекло;\\
    2. оно сцеплено с ключом держателя, привязанным эмитентом к удостоверению;\\
    3. именно это действие лежит внутри делегированной области;\\
    4. оно не отозвано, а удостоверение не тронуто;\\
    5. каждый отказ назван отдельно.};

  % --- the residue ---------------------------------------------------------------
  \node[futbox, anchor=north west, font=\scriptsize, text width=7.6cm, align=left]
    (res) at (8.4,-2.05) {
    \textbf{Не закрыто.} При обыкновенной привязке держателя служба видит ещё и
    значение токена удостоверения, поэтому вердикт называет соотносимость \emph{открыто}.
    Лишь путь с нулевым разглашением есть \emph{ограничено}.};

  % --- the property the whole design exists for ----------------------------------
  \node[anchor=north, font=\scriptsize\itshape, text=navymid, text width=15.5cm, align=center]
    at (7.5,-5.35) {Отзыв полномочия не отзывает человека. Держатель подписывает
    отзыв тем же ключом; к эмитенту никто не обращается, и он никогда не узнаёт, что полномочие
    существовало.};
\end{tikzpicture}
